\documentclass[11pt,a4paper]{article}
\usepackage[T1]{fontenc}
\usepackage[utf8]{inputenc}
\usepackage[french]{babel}
\usepackage{lmodern}
\usepackage[margin=2cm]{geometry}
\usepackage{microtype}
\usepackage{enumitem}
\usepackage{tabularx}
\usepackage{titlesec}
\usepackage{xcolor}
\usepackage[colorlinks=true,urlcolor=teal,linkcolor=black,pdfauthor={Charlotte Giraud},pdftitle={Curriculum vitae académique de Charlotte Giraud}]{hyperref}
\definecolor{teal}{HTML}{326D70}
\titleformat{\section}{\large\bfseries\color{teal}}{}{0pt}{}
\titlespacing*{\section}{0pt}{12pt}{6pt}
\setlength{\parindent}{0pt}
\setlength{\parskip}{4pt}
\setlist[itemize]{leftmargin=1.2em,itemsep=2pt,topsep=3pt}
\newcommand{\entry}[3]{\noindent\begin{tabularx}{\textwidth}{@{}>{\raggedright\arraybackslash}p{2.65cm}X@{}}\textit{#1}&\textbf{#2}\\& #3\end{tabularx}\par\vspace{3pt}}
\begin{document}
\begin{center}
{\LARGE\bfseries Charlotte Giraud}\\[5pt]
Enseignante certifiée · Historienne\\
Histoire économique et sociale · Genre\\[4pt]
\href{mailto:charlotte.giraudm@gmail.com}{charlotte.giraudm@gmail.com}
\end{center}

\section{Projet doctoral}
\textbf{Mobiliser ses possessions : genre, liquidité, réputation et crédit secret au Mont-de-Piété de Paris (1866--1918).}

Projet en préparation pour l’année universitaire 2027--2028. Corpus : 3\,921 engagements secrets dépouillés (1866--1877) et 6\,000 engagements secrets à dépouiller (1877--1918). Approche sérielle, quantitative et prosopographique, articulée à l’histoire du genre et des institutions de crédit.

\section{Activité scientifique}
\entry{24--28 juillet 2028}{World Economic History Congress (WEHC), Montevideo}{Contribution invitée au double panel « Financial Help in Hard Times: Credit, Assistance, and Social Protection in Historical Perspective » (panel en cours d’évaluation). Communication : « Secret Pawning and Financial Help: Reputation, Gender, and Social Protection at the Mont-de-Piété of Paris, 1866--1877 ».}
\entry{21--24 avril 2027}{European Social Science History Conference (ESSHC), Lyon}{Communication invitée : « Women as clients and creditors: financial networks and gender dynamics at the Mont-de-Piété de Paris, 19th Century ».}
\entry{11--12 septembre 2026}{Invisible Finance: Women and Money in History, Milan}{Communication : « Pawns, Pearls, and Power: Gendered Credit Practices at the Mont-de-Piété de Paris, 1866--1877 ».}
\entry{12--13 mars 2026}{Le genre et l’argent, le genre de l’argent et l’argent du genre, Lorient}{Colloque international interdisciplinaire, Université Bretagne Sud (HCTI / TEMOS). Publication associée : « Mobiliser ses possessions : minoration juridique et agentivité monétaire des femmes en France dans la deuxième moitié du XIXe siècle », avec Patrice Baubeau, à paraître dans \textit{Temps, Mondes et Sociétés} (début 2027).}

\section{Expérience de recherche}
\entry{Mars 2023--août 2026}{Assistante de recherche, Università Bocconi}{Projet EDIPO (\textit{Education, Diversity, Innovation and Politics}, 2022--2027). Recherche sur les élèves jésuites en France de la fin du XIXe au début du XXe siècle. Direction : Mara Squicciarini (Bocconi) et Nico Voigtländer (UCLA). Travail à distance au sein d’une équipe internationale : dépouillement, transcription, bibliographie et catalogage de sources.}
\entry{Juillet 2024--août 2025}{Assistante de recherche, University of California, Los Angeles (UCLA)}{Recherche sur la fiscalité sous Colbert, sous la direction de Nico Voigtländer (UCLA) et Benjamin Marx (Boston University).}

\newpage
\section{Formation}
\entry{2020--2023}{Université Paris Nanterre}{Master \textit{Histoire et sciences sociales du politique}, sous la direction de Patrice Baubeau et Laurence Croq, mention Bien. Mémoire : « Ma tante a un secret : elle préfère les diamants de la folie aux guenilles de la misère. Les pratiques irrégulières du prêt sur gage dans le registre d’engagements secrets du Mont-de-Piété de Paris à la fin du XIXe siècle » (173 p.). Bourse de recherche de la Mission historique de la Banque de France. Formation aux SIG à l’EHESS (2023).}
\entry{2019--2020}{Université Paris Nanterre}{Double licence Histoire-Anglais, mention Bien.}
\entry{2017--2019}{Lycée Gabriel Guist’hau, Nantes}{CPGE littéraire : hypokhâgne et khâgne A/L, spécialité Histoire-Géographie.}
\entry{2016--2017}{Lycée Joachim du Bellay, Angers}{Baccalauréat scientifique, spécialité Physique-Chimie, mention Très bien ; section européenne anglais.}

\section{Enseignement}
\entry{2026--2027}{Collège Jacques Prévert, Entrelacs}{Enseignante certifiée. Histoire-Géographie-EMC, classes de 5e et de 4e, Académie de Grenoble.}
\entry{2025--2026}{Lycée Pierre Béghin, Moirans}{Enseignante stagiaire. Histoire-Géographie-EMC, classes de seconde et de première, Académie de Grenoble.}
\entry{2020--2024}{Éducation nationale}{Assistante d’éducation. Lycée professionnel Claude Garamont (Colombes), Cité scolaire Molière et lycée professionnel Octave Feuillet (Paris 16e).}

\section{Méthodes}
Registres sériels ; critique des sources ; bases historiques ; statistiques descriptives ; prosopographie ; systèmes d’information géographique (SIG).

\section{Langues}
\begin{itemize}
\item Français : langue maternelle.
\item Anglais : courant, langue de travail.
\item Espagnol : intermédiaire.
\item Italien : débutant.
\end{itemize}

\section{Domaines de recherche}
Genre ; histoire économique et sociale ; crédit et dette ; prêt sur gage ; réputation et secret ; XIXe siècle.
\end{document}
